\section*{Bab \ref{ch:b1:mammal-organization} --- Organisasi Fungsional Seekor Mamalia}

\begin{solution}{exo:b1:mammal-organization:1}
Vertebrata: simetri bilateral, \omterm{def:b1:mammal-organization:bodyplan}{sefalisasi}, sumbu ruas tulang belakang
di punggung yang menyelubungi sumsum tulang belakang, \omterm{def:b1:mammal-organization:bodyplan}{selom} di sisi
perut berisi \omterm{def:b1:mammal-organization:organ}{organ} dalam, empat tungkai. Mamalia: \omterm{def:b1:mammal-organization:bodyplan}{diafragma} yang membagi
\omterm{def:b1:mammal-organization:bodyplan}{selom} menjadi dada dan perut, rambut, susu, suhu tubuh yang diatur.
\end{solution}

\begin{solution}{exo:b1:mammal-organization:2}
Nutrisi: pencernaan (lambung), pernapasan (paru), peredaran (jantung),
ekskresi (ginjal). Relasi: saraf (otak), indra (mata), endokrin
(tiroid), rangka dan otot (tulang paha, bisep), kulit. Reproduksi:
kelamin (ovarium, testis), kelenjar susu.
\end{solution}

\begin{solution}{exo:b1:mammal-organization:3}
Air $0.6\times 60 = \qty{36}{L}$; intrasel \qty{24}{L}; ekstrasel
\qty{12}{L}, yang antarselnya \qty{9}{L} dan plasmanya
\qty{3}{L}.
\end{solution}

\begin{solution}{exo:b1:mammal-organization:4}
$1.1/5 = 22\%$. \qty{1.1}{L/min} berarti \qty{1580}{L} sehari, yakni
volume darahnya 290 kali.
\end{solution}

\begin{solution}{exo:b1:mammal-organization:5}
Inulin yang tersisa \qty{1.7}{g}; $V = 1.7/0.12 = \qty{14}{L}$
ekstrasel; intrasel $42 - 14 = \qty{28}{L}$.
\end{solution}

\begin{solution}{exo:b1:mammal-organization:6}
Dinding kapiler meloloskan air dan ion kecil dengan bebas, sehingga
keduanya menyetimbang, tetapi menahan protein. Protein plasma menarik
air ke dalam pembuluh secara osmotik; tanpanya air akan merembes ke
ruang antarsel, yang akan membengkak (udem), dan volume plasmanya akan
turun.
\end{solution}

\begin{solution}{exo:b1:mammal-organization:7}
Secara paralel tiap \omterm{def:b1:mammal-organization:organ}{organ} menerima darah dengan oksigen arteri dan
tekanan penuh, dan alirannya dapat diatur sendiri-sendiri. Secara seri
otak akan menerima darah yang sudah dikuras otot, pada tekanan lebih
rendah, dan pasokannya justru akan turun selama olahraga.
\end{solution}

\begin{solution}{exo:b1:mammal-organization:8}
$hS\Delta T = 80 - 20 = \qty{60}{W}$; $h = 60/(1.8\times 13) =
\qty{2.6}{W.m^{-2}.K^{-1}}$. Pada udara \qty{33}{\celsius} kehilangan
bukan penguapan lenyap; kulitnya harus menghangat di atas suhu udara
(pelebaran pembuluh menaikkannya ke arah \qty{36}{\celsius}) dan
berkeringat harus membawa sisanya.
\end{solution}

\begin{solution}{exo:b1:mammal-organization:9}
Kalor yang harus dibuang $800 - 200 - 150 = \qty{450}{W}$, yakni
\qty{1620}{kJ/h}, \qty{675}{g} keringat tiap jam. Tanpa berkeringat,
$450/(70\times 3500) =
\qty{1.8e-3}{K/s}$, sekitar \qty{0.11}{K/min}: \qty{3}{K} dalam setengah jam.
\end{solution}

\begin{solution}{exo:b1:mammal-organization:10}
Tubuh kecil kehilangan kalor paling cepat tiap gram ($S/V$ besar) dan
memerlukan produksi yang sinambung, tinggi dan andal: lemak cokelat
menghasilkan kalor secara kimia pada laju tetap tanpa gerakan,
sedangkan menggigil menuntut otot yang sudah berkembang (yang belum
dimiliki bayi), mengganggu gerak dan makan, serta bersifat
terputus-putus. Lemak cokelat juga terletak dekat pembuluh besar dan
menghangatkan darah secara langsung.
\end{solution}

\begin{solution}{exo:b1:mammal-organization:11}
Hipotalamus yang dihangatkan membaca ``terlalu panas'': anjing itu
terengah-engah dan melebarkan pembuluh kulitnya meskipun dingin, dan
suhu intinya turun satu dua derajat dalam satu jam. Tanpa hipotalamus
tidak ada \omterm{def:b1:organism-environment:homeostasis}{titik acuan}: anjing itu tidak menggigil dalam dingin dan tidak
terengah dalam panas, dan suhunya melayang ke arah suhu ruang pada kedua
keadaan.
\end{solution}

\begin{solution}{exo:b1:mammal-organization:12}
\omterm{def:b1:mammal-organization:compartments}{Cairan ekstrasel} (\omterm{def:b1:mammal-organization:compartments}{cairan antarsel} dan plasma), yang susunan kaya
natriumnya, suhunya dan volumenya hampir tetap, merendam tiap sel
sebagaimana air laut merendam hewan pertama. Ketetapannya dijaga
gelung \omterm{def:b1:organism-environment:homeostasis}{umpan balik negatif} (ginjal untuk volume dan ion, paru untuk gas,
hati dan pankreas untuk glukosa, hipotalamus untuk suhu). ``Pantainya''
adalah \omterm{prop:b1:organism-environment:surfaces}{permukaan pertukaran} --- usus, paru, ginjal, kulit --- tempat ia
bertemu dunia luar, dan peredaran darah adalah arus yang mengaduknya.
\end{solution}

\begin{solution}{pb:b1:mammal-organization:1}
\textbf{1.} $B = 3.4\times 2^{0.75} = \qty{5.7}{W}$.
\textbf{2.} Liang: $5.7\times 12\times 3600 = \qty{246}{kJ}$; mencari
makan: $17.1\times 6\times 3600 = \qty{369}{kJ}$; di luar:
$11.4\times 6\times 3600 = \qty{246}{kJ}$.
\textbf{3.} \qty{861}{kJ}, yakni $861/493 = 1.75$ kali basal
($B$ selama \qty{24}{h} $= \qty{493}{kJ}$).
\textbf{4.} $0.25\times 18\times 0.65 = \qty{2.9}{kJ/g}$.
\textbf{5.} $861/2.9 = \qty{295}{g}$ rumput segar, \qty{74}{g} bahan
kering.
\textbf{6.} $295/2000 = 15\%$ massanya.
\textbf{7.} $861/20 = \qty{43}{L}$ $\mathrm{O_2}$; \qty{30}{mL/min}.
\textbf{8.} $0.9\times 43 = \qty{39}{L}$ $\mathrm{CO_2}$, yakni $39/22.4\times 44 = \qty{76}{g}$.
\textbf{9.} $0.75\times 295 = \qty{221}{g}$.
\textbf{10.} Bahan kering tercerna $0.65\times 74 = \qty{48}{g}$; air
$0.6\times 48 = \qty{29}{g}$.
\textbf{11.} Tak tercerna \qty{26}{g}; tinja $26/0.6 = \qty{43}{g}$,
yang membawa pergi \qty{17}{g} air.
\textbf{12.} $(44 - 5)\times 0.86 = \qty{34}{g}$.
\textbf{13.} Masuk: $221 + 29 = \qty{250}{g}$. Keluar: $17 + 130 + 34 + 20
= \qty{201}{g}$. Surplus \qty{49}{g}.
\textbf{14.} Tidak: rumputnya memasok lebih banyak daripada yang
hilang; surplusnya pergi sebagai kemih tambahan (sekitar \qty{180}{mL}
seluruhnya).
\textbf{15.} Bahan kering yang sama \qty{74}{g} dalam $\qty{148}{g}$
rumput: air dalam makanan \qty{74}{g}; pemasukan \qty{103}{g}
berhadapan dengan pengeluaran \qty{201}{g}: defisit \qty{100}{g} sehari,
yang harus diminumnya atau hilang dari tubuhnya.
\textbf{16.} Air $0.6\times 2.0 = \qty{1.2}{L}$; ekstrasel
\qty{0.4}{L}; plasma \qty{0.1}{L}.
\textbf{17.} $1.5/15 = \qty{0.10}{L}$: cocok.
\textbf{18.} $0.10/0.60 = \qty{0.167}{L}$, \qty{8}{\%} massa tubuh.
\textbf{19.} $5\times(2/70)^{0.75} = \qty{0.35}{L/min}$; waktu peredaran
$0.167/0.35 = \qty{0.48}{min}$, sekitar \qty{29}{s}.
\textbf{20.} $350/200 = \qty{1.75}{mL}$ tiap denyut.
\textbf{21.} $0.1\times 2^{2/3} = \qty{0.16}{m^2}$.
\textbf{22.} $2.5\times 0.16\times 34 = \qty{13.6}{W}$, berhadapan
dengan $2B = \qty{11.4}{W}$ yang diandaikan: si kelinci harus
menghasilkan sedikit lebih banyak, atau mengurangi kehilangan lewat
sikap tubuh dan tempat berlindung.
\textbf{23.} Terbuka: $10\times 0.01\times 34 = \qty{3.4}{W}$,
seperempat kehilangan lewat bulu lagi; menutup: $10\times 0.01\times 3 =
\qty{0.3}{W}$. Ia menutupnya.
\textbf{24.} Bulu: $2.5\times 0.16\times 9 = \qty{3.6}{W}$; telinga
terbuka: $10\times 0.01\times 9 = \qty{0.9}{W}$; seluruhnya \qty{4.5}{W}
berhadapan dengan \qty{5.7}{W} yang dihasilkan. Sisa \qty{1.2}{W} harus
pergi lewat penguapan: si kelinci terengah-engah (\qty{1.2}{W} setara
\qty{1.8}{g} air tiap jam), merebahkan diri di tanah yang sejuk dan
tinggal di tempat teduh.
\textbf{25.} Pengeluaran energi harian \qty{0.86}{MJ}, $1.75$ kali laju
basalnya; perputaran air \qty{250}{g}, \qty{12.5}{\%} massanya.
\end{solution}
